\section{Kapitel 4}
\subsection{Bibeltext}


    \spa{1}\hh{1}Nachdem also Christus für uns am Fleisch litt, \\
        \spa{2}wappnet auch ihr euch mit der selben Denkweise \\
            \spa{3}(weil der, der im Fleisch litt, mit Sünde abgeschlossen hat),\\
            \spa{3}\hh{2}um die noch verbleibende Zeit \\
            \spa{3}im Fleisch nicht mehr den Lüsten der Menschen zu leben, \\
                sondern dem Willen Gottes;\\
            \spa{3}\hh{3}denn es ist uns genug, \\
                \spa{4}die vergangene Zeit dieses Lebens \\
                \spa{4}den Willen derer, \\
                \spa{4}die von den Völkern sind, \\
                \spa{4}ausgeführt zu haben, \\
                \spa{4}als man sich erging \\
                    \spa{5}in Ausschweifungen, \\
                    \spa{5}Lüsten, \\
                    \spa{5}übermäßigem Trinken von Wein, \\
                    \spa{5}Schlemmereien, \\
                    \spa{5}Trinkgelagen \\
                    \spa{5}und sittenlosen Götzendiensten,\\
            \spa{3}\hh{4}wobei es sie befremdet, \\
                \spa{4}dass ihr nicht [mit ihnen] zusammen \\
                \spa{4}in denselben Erguss eines heillosen Wesens lauft, \\
            \spa{3}[sodass] sie lästern,\\
                \spa{4}\hh{5}[sie], \\
                \spa{4}die dem Rechenschaft geben werden, \\
                \spa{4}der in Bereitschaft steht, \\
                \spa{4}Lebende und Tote zu richten;\\
            \spa{3}\hh{6}denn zu diesem [Zweck] auch \\
                \spa{4}wurde [den] Toten gute Botschaft gesagt, \\
                \spa{4}dass sie zwar den Menschen gemäß am Fleisch gerichtet würden, \\
                \spa{4}aber Gott gemäß im Geist leben sollten. \\
    \spa{1}\hh{7}Aber das Ende aller Dinge ist nahegekommen. \\
        \spa{2}Seid also gesunden Sinnes \\
            \spa{3}und züchtig \\
            \spa{3}und seid nüchtern für die Gebete.\\
        \spa{2}\hh{8}Dabei habt vor allem zu euch untereinander \\
            \spa{3}eine spannkräftige Liebe \\
                \spa{4}(weil Liebe eine Menge von Sünden verdecken wird)\\
            \spa{3}\hh{9}als solche, \\
                \spa{4}die gastfreundlich gegeneinander sind ohne Murren \\
                \spa{4}\hh{10}und wobei [ihr], \\
                \spa{4}jeder so, \\
                \spa{4}wie er eine Gnadengabe empfing, \\
                \spa{4}euch damit dient \\
                \spa{4}– wie edle Hausverwalter der mannigfaltigen Gnade Gottes: \\
                    \spa{5}\hh{11}wenn jemand redet, [tue er es] als [einer, \\
                        \spa{6}der] Worte Gottes [spricht], \\
                    \spa{5}wenn jemand dient, \\
                        \spa{6}[sei es] als [einer, der es] aus der Stärkung [tut], \\
                        \spa{6}die Gott darreicht, \\
                    \spa{5}damit in allem Gott verherrlicht werde \\
                    \spa{5}durch Jesus Christus, \\
                        \spa{6}dem in alle Ewigkeit \\
                        \spa{6}die Herrlichkeit \\
                        \spa{6}und die Macht gebühren. \\
    \spa{1}Amen. \\
\\
    \spa{1}\hh{12}Geliebte, \\
        \spa{2}lasst euch den Brand, \\
            \spa{3}der unter euch zur Prüfung entstanden ist, \\
            \spa{3}nicht befremden, \\
            \spa{3}als widerführe euch etwas Fremdartiges, \\
        \spa{2}\hh{13}sondern so, \\
            \spa{3}wie ihr der Leiden des Christus teilhaftig seid, \\
            \spa{3}freut euch, \\
            \spa{3}damit ihr auch in der Enthüllung \\
                \spa{4}seiner Herrlichkeit euch jubelnd freuen mögt. \\
\\
        \spa{2}\hh{14}Wenn ihr im Namen des Christus, \\
        \spa{2}des Gesalbten, geschmäht werdet, \\
        \spa{2}[seid ihr] Selige, \\
            \spa{3}weil der Geist \\
            \spa{3}– der [Geist] der Herrlichkeit und Gottes – \\
            \spa{3}auf euch ruht; \\
                \spa{4}bei ihnen wird er gelästert, \\
                \spa{4}bei euch aber verherrlicht; \\
            \spa{3}\hh{15}es leide nämlich keiner von euch\\
                 \spa{4}als Mörder \\
                 \spa{4}oder Dieb \\
                 \spa{4}oder Übeltäter \\
                 \spa{4}oder als ein Aufseher fremder Angelegenheiten. \\
            \spa{3}\hh{16}Wenn [er] aber als Christ [leidet], s\\
                \spa{4}chäme er sich nicht, \\
                \spa{4}verherrliche aber Gott in diesem, \\
                \spa{4}das ihm zuteilwurde, \\
                    \spa{5}\hh{17}weil der Zeitpunkt [da ist], \\
                    \spa{5}an dem Gericht beginnen sollte, \\
                        \spa{6}[anfangend] beim Hause Gottes; \\
                    \spa{5}wenn aber zuerst bei uns, \\
                        \spa{6}was wird das Ende derer sein, \\
                        \spa{6}die der guten Botschaft Gottes im Unglauben ungehorsam sind? \\
            \spa{3}\hh{18}Und „wenn der Gerechte mit Mühe gerettet wird, \\
                \spa{4}wo wird der Ehrfurchtslose \\
                \spa{4}und Sünder erscheinen?“ {Spr 11,31 n. d. gr. Üsg.} \\
\\
            \spa{3}\hh{19}Daher sollen auch die, \\
                \spa{4}die nach dem Willen Gottes leiden, \\
                \spa{4}[ihm] als einem treuen Schöpfer \\
                \spa{4}ihre Seelen anvertrauen, [\\
                \spa{4}und dieses] im Gutestun. \\


